\documentclass[12pt]{article}

\newif\ifsol
\soltrue

% Unicode compatibility
\usepackage{iftex}
\ifPDFTeX
  \usepackage[utf8]{inputenc}
  \usepackage[noTeX]{mmap}
  \usepackage[T1]{fontenc}
\fi
% XeTeX does not support mmap
\ifLuaTeX
  \usepackage{luatex85}
  \usepackage[noTeX]{mmap}
\fi

\setlength{\textheight}{9in}
\setlength{\textwidth}{7.1in}
\setlength{\evensidemargin}{-0.2in}
\setlength{\oddsidemargin}{-0.2in}
\setlength{\headsep}{30pt}
\setlength{\topmargin}{-0.3in}

\usepackage{amsthm,amsfonts,amsmath,xspace,tikz,varwidth,cancel,soul}
\usepackage{hyperref,verbatim}
\newtheorem{theorem}{Theorem}
\theoremstyle{definition}
\newtheorem{definition}{Definition}

\newcommand{\addmedskip}{\addvspace{\medskipamount}}

\newcommand{\addbigskip}{\addvspace{\bigskipamount}}

\pagestyle{myheadings}
\markboth{BU CAS CS 538.}{BU CAS CS 538.} 
\newcounter{problemnum}
\setcounter{problemnum}{0}
\newenvironment{problem}
     {\addbigskip \setcounter{partnum}{0}
      \noindent\stepcounter{problemnum}\textbf{Problem
                                             \arabic{problemnum}.\ }}
     {\par\addbigskip}
     
\ifsol
\newenvironment{solution}
     {\addbigskip
      \noindent\color{blue}\textbf{Solution.\ }}
     {\par\addbigskip}
\else
\newenvironment{solution}
     {\comment}
     {\endcomment}
\fi

\newcounter{partnum}
\setcounter{partnum}{0}
\newenvironment{ppart}
     {\addmedskip
      \noindent\stepcounter{partnum}\textbf{(\alph{partnum})}\ }
     {\par\addbigskip}

\newcommand{\Enc}{\mathrm{Enc}}
\newcommand{\Dec}{\mathrm{Dec}}
\newcommand{\eqdef}{\buildrel{\scriptstyle{\mathit{def}}}\over{=}}
\newcommand{\rfrom}{\buildrel{\scriptstyle{\mathrm{R}}}\over{\leftarrow}}

\newcommand{\mbmod}{\,\%\,}

\newcommand{\A}{\mathcal{A}}
\renewcommand{\L}{\mathcal{L}}
\newcommand{\Z}{\mathbb{Z}}
\newcommand{\G}{\mathbb{G}}
\newcommand{\X}{\mathcal{X}}
\newcommand{\Y}{\mathcal{Y}}
\newcommand{\K}{\mathcal{K}}
\newcommand{\M}{\mathcal{M}}
\newcommand{\E}{\mathcal{E}}
\newcommand{\C}{\mathcal{C}}
\newcommand{\R}{\mathcal{R}}
\renewcommand{\S}{\mathcal{S}}
\newcommand{\I}{\mathcal{I}}
\newcommand{\T}{\mathcal{T}}
\newcommand{\randk}{\mathsf k}
\newcommand{\randm}{\mathsf m}
\newcommand{\randc}{\mathsf c}
\newcommand{\Adv}{\mathcal{A}}
\newcommand{\B}{\mathcal{B}}
\newcommand{\SSadv}{\textrm{SS}\textsf{adv}}
\newcommand{\PRGadv}{\textrm{PRG}\textsf{adv}}
\newcommand{\threebitadv}{\textrm{3Bit}\textsf{adv}}
\newcommand{\PRFadv}{\textrm{PRF}\textsf{adv}}
\newcommand{\MACadv}{\textrm{MAC}\textsf{adv}}
\newcommand{\CPAadv}{\textrm{CPA}\textsf{adv}}
\newcommand{\BCadv}{\textrm{BC}\textsf{adv}}
\newcommand{\CRadv}{\textrm{CR}\textsf{adv}}
\newcommand{\DDHadv}{\textrm{DDH}\textsf{adv}}
\newcommand{\comadv}{\textrm{Com}\textsf{adv}}
\newcommand{\pk}{\mathit{pk}}
\newcommand{\sk}{\mathit{sk}}
\newcommand{\qrn}{\mathit{QR}_n}
\newcommand{\qrp}{\mathit{QR}_p}
\newcommand{\qrpo}{\mathit{QR}_{p_1}}
\newcommand{\qrpt}{\mathit{QR}_{p_2}}
\newcommand{\crt}{\mathrm{crt}}




\newcommand{\indist}{  \mathrel{\vcenter{\offinterlineskip
  \hbox{$\sim$}\vskip-.35ex\hbox{$\sim$}\vskip-.35ex\hbox{$\sim$}}}}

\newcommand{\samples}{\xleftarrow{R}}

\newcommand{\hdl}{H_{\mathrm{dl}}}


\begin{document}
\begin{center}
\Large{\textbf{CAS CS 538.  \ifsol Solutions to \fi Problem Set 10}}\\
\smallskip
\large{\textbf{Due electronically via Gradescope, Wednesday April 22, 2026 11:59pm}}\\
\large{\textbf{Om Khadka, U51801771}}
\end{center}

\noindent
\textbf{How to submit:}  See instructions on Problem Set 1. 
 
\section{Commitments}


\begin{problem}[El Gamal commitments]
    We now consider a different commitment than the one from Discussion 10. 
    Let $\G$ again be a group of prime order $q$ generated by $g$. 
    Consider the following protocol between Alice and Bob. 
    Given a message $m\in \G$, Alice runs the following algorithm:
    \begin{itemize}
        \item \underline{\smash{Commit$(m)$}}: pick random $u\samples \G$,  $r\samples \Z_q$ and let $c=(u, g^r, u^r \cdot m)$.
    \end{itemize}
    Alice sends $c$ to Bob and keeps $r$ to herself.
    Bob makes sure that he got three elements of $\G$.
    
    Some time later, Bob wants to know what $m$ is and Alice is ready to reveal it.
    Alice sends $r$ and $m$ to Bob.
    Of course, she may not be honest, so we will call what Bob receives $m'$ and $r'$.
    Bob runs the following algorithm:
    \begin{itemize}
        \item \underline{\smash{Verify$(c, m',r')$}}: 
        First, make sure that $m' \in \G$ and $r' \in \Z_q$. 
        Let $c=(u, v, z)$. Check that $v=g^{r'}$ and that $z=u^{r'} m'$. 
        If so, accept $m'$ as the correct message. Else, reject.
    \end{itemize}
    Note that this commitment scheme is just the encryption scheme known as 
    multiplicative ElGamal (Exercise 11.5), but without decryption.\footnote{
    In fact, a similar commitment scheme can be built out of any encryption 
    scheme that uses randomness $r$: to commit to $m$, output $\pk$ and $E(\pk, m)$ 
    where $E$ uses randomness $r$ and no other randomness. To verify a revealed 
    message $m'$, use $r'$ to re-encrypt $m'$ under $\pk$ and check that the result 
    matches the received commitment. It works as long as for every public key --- 
    even adversarially generated --- correct decryption is unique.}

    \begin{ppart}
        Prove that this commitment scheme is \emph{computationally hiding} 
        under the DDH assumption.
        
        The definition of computationally hiding commitments is essentially 
        the same as the definition of semantically security encryption schemes: 
        given $\G, g$, the adversary outputs $m_0$ and $m_1$, 
        gets a commitment to $m_b$ for random bit $b$ and outputs $\hat{b}$, 
        with the usual measure of adversarial advantage 
        $\comadv = \lvert\Pr[\hat{b}=1\, |\,b=0]-\Pr[\hat{b}=1\,|\,b=1]\rvert$.
    \end{ppart}
    \begin{solution}
Your solution goes here
\end{solution}


\begin{ppart}
Prove that the commitment is \emph{perfectly binding}: namely, no matter how clever or computationally unbounded Alice is, for every $c$ and every $r_1, r_2$, there is at most one $m$ that Bob will accept.
  In other words, for every $c, r_1, r_2, m_1, m_2$, if Verify$(c, m_1, r_1)$ accepts and Verify$(c, m_2, r_2)$ also accepts, then $m_1 = m_2$.
\end{ppart}

\begin{solution}
Your solution goes here
\end{solution}
\end{problem}

\begin{problem} (20 points)
The above scheme is \emph{perfectly} binding and \emph{computationally} hiding. In discussion 10 we covered Pedersen commitments, which are \emph{computationally} biding and \emph{perfectly} hiding. 

To recap the discussion, a commitment scheme is \emph{perfectly hiding} when
for every $g, h, m_0$ and $m_1$, the distribution of outputs 
produced by Commit$(m_0)$ is identical to the distribution of outputs 
produced by Commit$(m_1)$. (Note that this protects Alice during the commit phase because Bob learns nothing about $m$, regardless of Bob's computational power or any computational assumptions.)

Prove that it's not possible to have a commitment scheme to be both perfectly binding and perfectly hiding.
\end{problem}
\begin{solution}
Your solution goes here
\end{solution}




\section{Weak PRF to PRF}

Recall that in the attack game for a PRF $F$ defined over $(\K, \X, \Y)$,
the adversary $\Adv$ submits a sequence of queries to the challenger, 
and the challenger picks a random key $k$ and answers each query $x_i$ with $F(k, x_i)$ in Experiment 0,
and it picks a random function $f$ from the set of all functions from $\X$ to $\Y$ and answers each query 
$x_i$ with $f(x_i)$ in Experiment 1.
Note that the adversary has full control over the queries.

We can, however, define a weaker notion of a PRF, called a \emph{weak PRF}, 
where the adversary is only allowed to make function queries on random inputs.
More specifically, in the weak PRF game, whenever the adversary $\Adv$ makes a query,
the \emph{challenger} picks a random $x_i \in \X$ and answers the query with $(x_i, F(k, x_i))$ in Experiment 0, 
and with $(x_i, f(x_i))$ in Experiment 1.
As expected, the adversary's advantage in this game, denoted $\mathrm{wPRF}\textsf{adv}[\Adv, F]$, 
is defined as the difference between the probability that $\Adv$ outputs 1 in Experiment 0 and the probability that $\Adv$ outputs 1 in Experiment 1; 
and the PRF $F$ is said to be a weak PRF if $\mathrm{wPRF}\textsf{adv}[\Adv, F]$ is negligible for all 
efficient adversaries $\Adv$. (See definition 4.3 in Section 4.4.1 of the textbook if you want more explanations.)

\begin{problem} (35 points)
Problem 11.1 in the textbook. Correction to the problem: let $q$ be the total number of queries that $\Adv$ makes. Then
$$\mathrm{PRF}^{\mathrm{ro}}\mathsf{adv}[\mathcal{A}, \mathcal{F} ] \le \mathrm{wPRF}\mathsf{adv}[\mathcal{B}, \mathcal{F'}]$$
should be instead
$$\mathrm{PRF}^{\mathrm{ro}}\mathsf{adv}[\mathcal{A}, \mathcal{F} ] \le \mathrm{wPRF}\mathsf{adv}[\mathcal{B}, \mathcal{F'}]+\frac{q^2}{2\cdot |\mathcal{X}|}$$

\paragraph{Interpretation:} We know that $\Adv$ is polytime bounded and hence it can make only a polynomial number of queries. Thus, $q^2$ is negligible if $|\mathcal{X}|$ is exponentially large. So this weak-PRF-to-PRF conversion works for large $|\mathcal{X}|$ (something like $256$-bit is reasonable) but not for small $\mathcal{X}$. 


\paragraph{Approach:}
	Let $\Adv$ be a PRF adversary against $F$. Because we are in the random oracle model, $\Adv$ gets to ask two types of queries: 
\begin{itemize}
	\item $q_p$ PRF queries $m_i$ (for $1\le i \le q_p$) to the PRF (in Game 0) or random function (in Game 1)
	\item $q_H$ random oracle (RO) queries $\hat{m}_i$ (for $1\le i\le q_H$) to $H$	
\end{itemize}
The total number of queries is $q=q_p+q_H$.
	
	Let $p_0$ be the probability that $\Adv$ outputs $1$ in experiment 0 of the PRF game (queries answered by $F$), and $p_1$ be the probability that $\Adv$ outputs $1$ in experiment 1 of the PRG game(queries answered by a random $f$). 
	
	Let $\C$ be the following challenger: pick a random function $H: \M \to\X$ and a random function $f':\X\to \Y$;	answer PRF queries $m_i$ with $f'(H(m_i))$ and RO queries $\hat{m}_i$ with $H(\hat{m}_i)$. Let $p_\C$ be the probability that $\Adv$ outputs $1$ when interacting with $\C$. 
	
	In the text two parts, you will prove $|p_0-p_c|$ and $|p_c-p_1|$ are both small, giving you that $|p_0-p_1|$ is small.

	\begin{ppart}
	Prove that $|p_0-p_c|\le \mathrm{wPRF}\mathsf{adv}[\mathcal{B}, \mathcal{F'}]$ by building an adversary $\B$ that uses $\Adv$ and plays the wPRF game of $F'$ with advantage  $|p_0-p_c|$. 
	
	Hint: $\B$ has to figure out how to consistently answer RO queries of $\Adv$. For example, if $\Adv$ asks $\B$ for $H(m)$, $\Adv$ may later ask $\B$ for $F'(H(m))$ (or $f'(H(m))$, depending on which game it's in). $\B$ should prepare for this eventuality, because it can't ask its wPRF challenger for a particular input. (This is called ``preparing'' or ``programming'' the random oracle.)	
	\end{ppart}
	
	\begin{solution}
Your solution goes here
\end{solution}


\begin{ppart}
Prove that $|p_c-p_1|\le \frac{q^2}{2\cdot|\X|}$ by the difference lemma.
\end{ppart}	
	

\begin{solution}
Your solution goes here
\end{solution}
	
\end{problem}


\end{document}